SafeEdit-CRE addresses a central challenge in computational regulatory genomics: how
to design minimal, cell-type-selective edits to natural CREs such that the resulting
variants generalize across sequence-to-activity models rather than exploiting a single
predictor. Across three complementary benchmarks---600 held-out natural parents, a
96-parent nine-criterion design-quality benchmark, and seven complete re-search
ablations---SafeEdit-CRE produced three principal findings.

First, relative to iterative greedy optimization, SafeEdit-CRE improved cross-model
cell-type-specificity margin by 0.055 standardized units (95\% CI 0.040--0.071).
This is a meaningful gain in a setting where the greedy baseline already achieves
positive margin gain in 97.1\% of conditions, and it demonstrates that single-model
local optimization leaves substantial cross-model transfer performance on the table.
The improvement was most pronounced at edit budgets of 5--20 nucleotides in K562 and
HepG2, where the larger space of plausible edit paths increases both the opportunity
for overfitting and the value of independent review.

Second, against a compute-matched primary-model beam that explores the same number of
candidate paths but ranks them solely by primary-predictor score, SafeEdit-CRE
preserved predicted specificity within tight equivalence bounds (difference
$-0.004$; 95\% CI $-0.011$--0.004) while lowering cross-model ensemble uncertainty by
0.021 (95\% CI $-0.025$ to $-0.017$) and reducing aggressive target-gain excursions.
This equivalence-and-improvement pattern is a favorable design outcome: the
cell-type-specific objective is fully retained, while the sequences returned are less
dependent on any single model's idiosyncratic extrapolations. In settings where
oligonucleotide synthesis and reporter assays are the limiting resources, preserving
predicted selectivity while reducing predictive dispersion is more useful than simply
maximizing the largest score available from one model.

Third, SafeEdit-CRE raised nine-criterion design compliance from 38.5\% (greedy) to
61.5\%, a 22.9 percentage-point improvement that reflects the value of enforcing
sequence-domain and cross-model constraints \emph{during} search rather than filtering
afterward. The complete re-search ablation confirmed that naturalness-aware reranking
is the dominant contributor to uncertainty reduction, while the sequence-domain
prefilter enforces synthesis feasibility (GC and homopolymer limits) at a modest cost
in predicted margin. The reviewer ensemble is the component that distinguishes
SafeEdit-CRE from the primary beam: removing it reproduced the primary beam exactly.

The tiered candidate library changes the practical form of the experimental output.
Instead of returning a single highest-scoring sequence per parent, SafeEdit-CRE
produces candidate-level records for primary activity, reviewer transfer, ensemble
uncertainty, strand consistency, naturalness, base composition, homopolymers, and edit
count. The 78 Tier A sequences satisfy all nine criteria and lie on or near the
primary--reviewer Pareto frontier, spanning all three cell types and four budgets to
provide a compact first-pass synthesis panel. The additional 1,224 Tier B alternatives
retain breadth for follow-up screening. Three representative Tier A examples
(Supplementary Table~S4) illustrate the range of edits, from single substitutions
that subtly shift predicted specificity to multi-base edits that produce larger
margin gains while maintaining sequence plausibility; observed motif-level changes are
presented as testable hypotheses rather than asserted mechanisms.

The study was deliberately built around evaluation practices that strengthen the
interpretation of model-guided design results~\citep{Whalen2022,deBoer2024,Rafi2025}.
Parent selection was label-blind; development and evaluation parents were separated by
identifier and sequence hash; random replicates were averaged rather than selected by
outcome; confidence intervals resampled whole parents to preserve correlation
structure; and the independent evaluation ensemble was excluded from candidate
generation. The compute-matched beam control separates the benefit of search breadth
from the benefit of cross-model reranking, a distinction that is absent from most
published sequence-design evaluations. The complete re-search ablation ensures that
component effects reflect changes in which sequences are selected, not retrospective
filtering of a fixed candidate set.

The present study establishes computational performance and candidate prioritization.
Prospective reporter assays will provide the next level of validation for cellular
activity. Importantly, the independent cross-model evaluation design, compute-matched
controls, and full re-search ablations provide a rigorous basis for selecting
candidates before committing experimental resources. The reviewer and evaluation
ensembles, while architecturally distinct from the primary predictor, share the same
MPRA training data, and the 6-mer naturalness score is a sequence-domain proxy rather
than a model of chromatin context; however, the cross-model disagreement captured by
these controls is precisely the signal that the method exploits to avoid overfitting,
and the experimental candidate panel is designed to test those predictions directly.

SafeEdit-CRE demonstrates that cross-model selection can improve the reliability of
minimal regulatory-DNA editing without sacrificing the predicted cell-type-specific
objective. Its combination of natural-sequence scaffolds, uncertainty-aware ranking,
and independent evaluation offers a practical route from sequence-to-activity models
to compact experimental candidate libraries. The same design principle can be applied
to other regulatory-prediction models and other cellular contexts whenever a limited
synthesis budget makes candidate reliability as important as the predicted optimum.
